% ================================================================
% CONTOH — PENGGUNAAN CLASS FILE
% Cukup tulis \documentclass{academic-id}; seluruh framework dimuat
% otomatis (metadata, packages, settings, command).
%
% Identitas dokumen diisi di preamble/metadata.tex.
% Kompilasi dari ROOT proyek (agar academic-id.cls ditemukan):
%   latexmk -lualatex -outdir=output examples/class/main.tex
% ================================================================

\documentclass{academic-id}

% Class memakai framework dari root; metadata (\docLogo, \docBibliography)
% diambil dari preamble/metadata.tex yang sudah relatif root.

\begin{document}

\makecover[][Proposal Penelitian]

\chapter*{Abstrak}
Penelitian deskriptif ini bertujuan untuk memberikan gambaran umum mengenai topik yang diteliti. Bagian ini dapat diisi dengan ringkasan masalah, metode, hasil, dan kesimpulan penelitian secara singkat.

\vspace{1em}
\textbf{Kata kunci:} kata kunci 1, kata kunci 2, kata kunci 3

\cleardoublepage


\chapter{Pendahuluan}
\section{Latar Belakang}
Contoh pemakaian class file \texttt{academic-id}. Semua pengaturan dan
perintah framework tersedia otomatis, tanpa perlu meng-import preamble
secara manual.\supercite{santosoKarakteristikGlaukomaTraumatik2020}
\section{Tujuan}
Menunjukkan cara memakai framework sebagai class file.

\chapter{Metode}
\section{Desain}
Contoh deskriptif sederhana.

\Table{tab:contoh}{
  \begin{tabular}{ll}
  \toprule
  Variabel & Keterangan \\
  \bottomrule
  \end{tabular}
}{Contoh tabel dengan perintah \texttt{\textbackslash Table}.}

\printbibliography[title={Daftar Pustaka}]

\end{document}
